%% Abstract
\chapter*{Abstract}
\addcontentsline{toc}{chapter}{Abstract}

This report presents a comprehensive biodiversity delta analysis for the site known as Pingle, Taylor's Lane, Ashleyhay, Derbyshire (DE4~4AH), covering a 1\,km radius study area centred on the property. The analysis employs the Statutory Biodiversity Metric~4.0 (BM\,4.0) framework, as mandated by the Environment Act 2021 and its associated regulations (Statutory Instruments 2024 No.~47), to quantify changes in habitat value between a 2016 baseline (\Tzero) and a 2026 present-day assessment (\Tone).

The \BDE{} pipeline integrates multi-source Earth observation data—including Sentinel-2 multispectral imagery, Landsat~8 OLI/TIRS composites, and Ordnance Survey MasterMap topography layers—with supervised machine learning classification (Random Forest ensemble) to produce parcel-level UK Habitat Classification assignments at both temporal epochs. These classifications are then scored against BM\,4.0 distinctiveness, condition, and strategic significance matrices to derive area-based Biodiversity Units (BU) for each parcel.

Key findings indicate a net delta of \dataplaceholder{NET DELTA BU} biodiversity units across the study area, representing a \dataplaceholder{PERCENTAGE}\% change from the 2016 baseline. The analysis identifies \dataplaceholder{N IMPROVED} parcels with improved habitat condition, \dataplaceholder{N DEGRADED} parcels showing degradation, and \dataplaceholder{N STABLE} parcels with negligible change. Hedgerow and watercourse linear features are assessed separately, yielding \dataplaceholder{HEDGEROW DELTA} linear hedgerow units and \dataplaceholder{WATERCOURSE DELTA} river units of change.

Monte Carlo uncertainty quantification across 10,000 iterations provides 95\% confidence intervals for all reported delta values. Sensitivity analysis demonstrates that the results are robust to reasonable variation in condition scoring assumptions. Cross-validation against published ecological survey data and a University of Derby undergraduate thesis on the same study area confirms alignment within expected tolerances.

This document serves as both a technical reference for the computational methodology and a statutory-grade evidence base suitable for submission to Derbyshire County Council in support of biodiversity net gain obligations.

\vspace{5mm}
\noindent\textbf{Keywords:} Biodiversity Net Gain, Statutory Biodiversity Metric~4.0, delta analysis, habitat classification, Derbyshire, Earth observation, machine learning, uncertainty quantification.
